\chapter{Background} % (fold)
\label{cha:background}

\begin{abstract}
In the previous chapter, we have defined \ac{tg} as the problem of finding a
sequence of actions that will move the robot from its current state to a
desired state. As this thesis focuses on \ac{tg} for mobile manipulation,
relevant literature must therefore include works from different fields, ranging
from autonomous driving and mobile robotics to manipulation. Besides, \ac{tg} in
dynamic environments generally includes path planning and local \ac{tg}.
Therefore, this chapter summarizes approaches ordered on a scale of reactivity,
that is the frequency at which trajectories are computed. Starting with
controller-like \ac{tg} methods, we then move in order of increasing
time-horizon up until global path planning.
\end{abstract}


\newpage

\section{Trajectory Generation} % (fold)
\label{sec:trajectory_generation}

In this thesis, we define \ac{tg} as the process of
computing a sequence of actions in real-time that allows a
robot to move from an initial configuration to a desired
goal state. Imporantly, the goal state might be defined in a
task specific way. For example, the goal state might be
defined in the task space (e.g., the end-effector position)
or in the configuration space (e.g., the joint position).
That places \ac{tg} between task planning and control.

To relate this definition of \ac{tg} to the literature,
\ac{tg} can be further split into global path planning and
local \ac{tg}. The former is the process of
computing a collision-free path from the initial to the
goal state, while the latter is the process of computing a
smooth trajectory that follows the path as closely as
possible while satisfying the robot's kinematic and dynamic
constraints ~\cite{Rickert2014}.

% section trajectory_generation (end)

\section{Dynamic Environments} % (fold)
\label{sec:dynamic_environments}

Modern research in robotics is increasingly focusing on
robots that operate in dynamic environments. However,
dynamic environments is a broad term that can refer to
environments with a variety of different characteristics. In
this thesis, we define a dynamic environment as an
environment in which the robot's task is, generally, not
known prior to deployement. That means that goal configurations, such as
goal locations, the size and shape of manipulation objects,
may vary during deployment. Besides, the environment may
contain moving obstacles, such as other robots, humans, or
other unpredictable objects. Examples of such environments
are supermarkets, hospitals, or homes, see \cref{fig:dynamic_environments}.

% Figure with two subfigures showing a hospital and a supermarket
\begin{figure}[ht]
  \centering
  \begin{subfigure}[b]{0.45\textwidth}
    \includegraphics[width=\textwidth]{robot_in_hospital.png}
  \end{subfigure}
  \hfill
  \begin{subfigure}[b]{0.45\textwidth}
    \includegraphics[width=\textwidth]{human_shared_environment.png}
  \end{subfigure}
  \caption{Examples of dynamic environments.}
  \label{fig:dynamic_environments}
\end{figure}


% section dynamic_environments (end)

\section{Notations} % (fold)
\label{sec:notations}

Throughout this thesis, vectors are denoted in bold lowercase
letters, $\vec{x}$, matrices in capital, $\mat{M}$, and sets
in calligraphic uppercase, $\mathcal{M}$. In the context of
discrete-time systems, we denote the time step using a
subscript, e.g., $\vx_k$ is the state at time $k$. The
transpose of a matrix is denoted by a superscript $T$, and
the inverse by a superscript $-1$, the pseudo-inverse by a
superscript $\dagger$. Partial derivates are denoted by 
$\der{\vx}{\vy}$ in text block, or more explicitely as
\[\derf{\vx}{\vy}.\]
 $\norm{\vx}_M = \vx^TM\vx$
denotes the weighted squared norm. 


We denote $\q\in\Q\subset\Rn$ a configuration of the robot
with $n$ its degrees of freedom; \Q{} is the configuration
space of the generalized coordinates of the system.
Generally, $\q(t)$ defines the robot's configuration at time
$t$, so that \qdot{}, \qddot{} define the instantaneous time
derivatives of the robot's configuration. Similarly, we
assume that there is a set of task variables
$\xj\in\Xj\subset\Rmj$ with variable dimension $m_j \leq n$.
The task manifold \Xj{} defines an arbitrary manifold of the
configuration space \Q{} in which a robotic task can be
represented. Further, we assume that there is a smooth differential
map $\map_j:\Rn\rightarrow\Rmj$ that relates the
configuration space to the $j^{th}$ task space. For example,
when a task variable is defined as the end-effector
position, then $\map_j$ is the positional part of the
forward kinematics. On the other hand, if a task variable is
defined to be the joint position, then $\map_j$ is the
identity function. In the following, we drop the subscript
$j$ in most cases for readability when the context is clear.

In this work, we assume that \map{} is twice
differentiable so that the Jacobian is defined as
\begin{equation}
  \J = \derf{\q}{\map} \in \mathcal{R}^{m\times n}, 
\end{equation}
or $\J = \der{\q}{\map}$ for short.
Thus, we can write the total time derivatives of \x{} as
\begin{align}
  \xdot &= \J\qdot \\
  \xddot &= \J\qddot + \Jdot\qdot.
\end{align}

% section notations (end)

\section{Model Predictive Control} % (fold)
\label{sec:model_predictive_control}

Optimization methods, such as dynamic programming are
populare when addressing \ac{tg} problems that are
characterized by multiple equality and inequality
constraints. 
The optimization problem is composed of a cost function, 
a dynamic model of the system, and a set of constraints.

This thesis uses exclusively discrete time dynamics, such
that the dynamics of the system are given by 
\begin{equation}
  \z_{k+1} = f(\z_k, \u_k),
\end{equation}
where $\z_k \in \Rn$, $\u_k \in \Rm$ are the state and
control input vectors, respectively. 

Multiple objectives of the \ac{tg} problem are combined in
the cost function:
\begin{equation}
  J(\z,\u) = \sum_{k=0}^{N-1} J_k(\z_k, \u_k) + J_f(\z_N,\u_N),
\end{equation}
where $N$ is the prediction horizon, and $J_k$ and $J_f$ are
the stage and terminal cost, respectively.

The general \ac{mpc} problem is then formulated as
\begin{equation}
  \begin{aligned}
    \min_{\u} \quad & J(\z,\u) \\
    \text{subject to} \quad & \z_{k+1} = f(\z_k, \u_k), \quad k=0,\ldots,N-1 \\
    & g(\z_k, \u_k) \leq 0, \quad k=0,\ldots,N-1 \\
    & \z_k \in \Z, \quad \u_k \in \U, \quad k=0,\ldots,N \\
    & \z_0 = \z_{\textrm{init}},
  \end{aligned}
\end{equation}
where $\z_0$ is the initial state, $\Z$ and $\U$ are the
admissibe sets of states and control inputs, respectively.
Inequality constraints are denoted by $g(\z_k, \u_k)$.

At each time step $k$, the \ac{mpc} problem is solved to
find the sequence of control inputs $\u_k^*$ that minimize
the cost function. The first control input of the optimal
sequence is then applied to the system. The remainder of the
sequence is shifted by one step to form the initial guess
for the next solution. This process is usually refered to as
receding horizon control.

\section{Differential Geometry} % (fold)
\label{sec:differential_geometry}

In this section, we consiously use a simplified formalism
for differential geometry to allow its understanding for
engineers. For a more thorough introduction to differential
geometry, we refer to ~\cite{bullo2019geometric}.

Differential geometry is the study of geometry using calculus.
It is a mathematical discipline that uses the techniques of
differential and integral calculus, as well as linear
algebra, to study problems in geometry. In the context of
this thesis, we recall some of the basic concepts of
differential geometry that are required to understand the
formalism of \ac{fabrics}.

\subsection{Manifolds} % (fold)
\label{sub:manifolds}

A manifold is a topological space that locally resembles
Euclidean space near each point. More precisely, a manifold
is a topological space $M$ such that for each point $p\in
M$, there exists a neighborhood $U$ of $p$ that is
homeomorphic to an open subset of $\Rn$. The dimension of
the manifold is the dimension of the Euclidean space to
which it is locally homeomorphic.

A smooth manifold is a manifold that is equipped with a
smooth structure. A smooth structure is a collection of
charts that cover the manifold and are compatible with each
other. A chart is a homeomorphism from an open subset of the
manifold to an open subset of $\Rn$. The compatibility
condition is that the transition maps between overlapping
charts are smooth. A smooth map between two smooth manifolds
is a map that is smooth in the sense that it is smooth in
each chart.

% subsection manifolds (end)

% section differential_geometry (end)

\section{Optimization Fabrics} % (fold)
\label{sec:optimization_fabrics}

In this section, we introduce the concept of \acf{fabrics}
\ac{fabrics} are a framework for \ac{tg}
that is based on the theory of differential
geometry. The main idea is to design \ac{tg} as a
second-order dynamical system. The trajectory generator is
defined by the differential equations as we will layout in
the following.

\subsection{Spectral Semi-Sprays} % (fold)
\label{sub:spectral_semi_sprays}

Let us first define the base class for all that follows.
Let $\X$ be a smooth manifold for which we denote vectors as 
$\x\in\X$ and $\xdot\in T\X$ the tangent space of $\X$.

Then, a \textit{spectral semi-spray}~\cite{Ratliff2020} is a
second-order dynamical of the form $\M\xddot + \f = 0$,
where $\M(\x,\xdot)$ and $\f(\x,\xdot)$ are functions of
position and velocity. Besides, \M{} is symmetric and
invertible. We often denote \textit{specs} using the shorter
notation $\Spec = \spec$, where we drop the subscript when
the space of the task variable is clear from the context.
Then the trajectory is defined as the solution to the system
$\xddot=-\M^{-1}\f$.

% subsection spectral_semi_sprays (end)

\subsection{Operations on Specs} % (fold)
\label{sub:operations_on_specs}

On \textit{specs}, we can define two fundamental operations:
the \textit{pullback} and the \textit{summation}.

Given two specs $\Spec_1 = {\left(\M_1,\f_1\right)}_{\X}$ and
$\Spec_2 = {\left(\M_2,\f_2\right)}_{\X}$, their sum is
defined as
\begin{equation}
  \Spec_1 + \Spec_2 = {\left(\M_1 + \M_2, \f_1 + \f_2\right)}_{\X}.
  \label{eq:summation}
\end{equation}

Given a differential map $\map: \Q\rightarrow\X$ and a spec
\spec{}, the \textit{pullback} is defined as
\begin{equation}
  \pull{\map}{\spec} = {\left(\Jt\M\J, \Jt(\f+\M\Jdot\qdot)\right)}_{Q}.
  \label{eq:pullback}
\end{equation}

The pullback allows converting between two distinct
manifolds (e.g. a spec could be defined in the robot's
workspace and being pulled into the robot's configuration
space using the pullback with \map{} being the forward
kinematics).

\subsection{Optimization fabrics}%
\label{sub:optimization_fabrics}

\ac{fabrics} form a special class of specs, and thus they inherit their properties,
specifically the previously defined operations of \textit{summation} and \textit{pullback}.
First, let us introduce a finite and differentiable potential
function $\forc(\x)$ defined in a task manifold \X{}. 
Then, the modified spec $\Spec_{\forc} = \left(\M,\f + \der{\x}{\forc}\right)$
is called the \textit{forced variant} of $\Spec = \spec$.
Only if the trajectory $\x(t)$ generated by the forced spec converges to the minimum of \forc{}, 
the spec is said to form an \textit{optimization fabric}.
When the spec only converges to the minimum when equipped with a damping term,
$\left(\M,\f + \der{\x}{\forc} + \mat{B}\xdot\right)$, 
it forms a \textit{frictionless fabric}~\cite[Definition 4.4]{Ratliff2020}. 
Note that the mechanical system of a pendulum forms a frictionless fabric, as it optimizes
the potential function defined by gravity when being damped (i.e., it eventually comes to
rest at the configuration with minimal potential energy)

In the following, methods to construct optimization fabrics, or \textit{fabrics} for
short, are summarized: the definitions of conservative fabrics and energization are
introduced.

\subsection{Conservative fabrics and energization}%
\label{sub:conservative_fabrics_and_energization}

While the previous subsection defined what criteria are required for a spec to form
an optimization fabric, the theory on conservative fabrics and energization 
offers a simple way of generating such special specs. As a full summary of the theory
on optimization fabrics and their construction is out of scope here, this
subsection only provides an outline of the theory and the reader is referred to 
\cite{Ratliff2021,Wyk2022} for detailed derivations.

In the context of \ac{fabrics}, the term \textit{energy} describes a scalar quantity that
changes as the system evolves over time.  Although this quantity has a physical meaning in
natural systems (e.g., kinetic energy), it can be arbitrarily defined for motion generation.
Besides, this quantity can be used as
a Lyaupnov function to prove the convergence of the trajectory to the minimum of the potential
function.
Generally, specs and optimization fabrics do not conserve an energy, but when they do, we
call them \textit{conservative specs}.  A stationary Lagrangian~\cite[Definition
4.11]{Ratliff2020} is one definition for an energy for which the corresponding spec, known
as the Lagrangian spec $\Spec_{\le} = \left(\Me,\fe\right)$, is obtained by applying the
Euler-Lagrange equations. Importantly, Lagrangian specs conserve energy and do thus belong
to the class of \textit{conservative specs}.  It was proven that an unbiased
(\cite[Definition~4.11]{Ratliff2020}) Lagrangian spec forms a frictionless
fabric~\cite[Proposition~4.18]{Ratliff2020}. Such fabrics are analogously called
\textit{conservative fabrics}.  There are two classes of conservative fabrics: Lagrangian
fabrics (i.e., the defining energy is a Lagrangian) and the more specific subclass of
Finsler fabrics (i.e., the
defining energy is a Finsler structure~\cite[Definition 5.4]{Ratliff2020}). 

The operation of \textit{energization} transforms a given differential equation into
a conservative spec.
Specifically, given an unbiased energy Lagrangian \le{} with boundary conforming
\Me{}~\cite[Definition~4.6]{Ratliff2020} and
lower bounded energy \he{}, an unbiased spec of form $\Spec_{\vec{h}} = (\mat{I},\vec{h})$
is transformed into a frictionless fabric using energization as
\begin{equation}
  \begin{split}
  S_{\vec{h}}^{\le} &= \text{energize}_{\le}\{S_{\vec{h}}\} \\
    &= (\Me, \fe + \Pe[\Me\vec{h} - \fe]), 
  \end{split}
\end{equation}
where $\Pe = \Me\left(\Me^{-1} - \frac{\xdot\xdot^T}{\xdot^T\Me\xdot}\right)$ is an
orthogonal projector.
Energized specs maintain the energy of the Lagrangian and generally change
the trajectory of the underlying spec $\Spec_{\vec{h}}$.
However, if 
\begin{enumerate}
  \item $\Spec_{\vec{h}} = (\mat{I},\vec{h})$ is homogeneous of degree 2,
    \begin{equation}\vec{h}(\x, \alpha\xdot) = \alpha^2\vec{h}(\x, \xdot)\label{eq:homogeneous}\end{equation}
    and
  \item the energizing Lagrangian is a Finsler structure, 
\end{enumerate}
the resulting energized spec forms a frictionless fabric for which the trajectory matches
the original trajectory of $\Spec_{\vec{h}}$. We refer to energized fabrics with that
property as \textit{geometric fabrics}. Geometric fabrics form the building blocks for
motion generation with optimization fabrics.
Practically, energization equips the individual components of the planning problem
with a metric when being combined with other components.


\subsection{Motion generation using optimization fabrics}%
\label{sub:trajectory_generation_using_optimization_fabrics}

After this brief introduction to the theory of \ac{fabrics},
we can layout the general procedure for \ac{tg} with the
framework.
When using optimization fabrics for \ac{tg}, all
components including constraints and goal attraction are
designed as second order differential equations. If specific
design rules for these equations are respected, all
components can be combined to form a converging motion
generator. Specifically, the following steps are performed:
%
\begin{enumerate}
  \item Design path-consistent geometries in suited manifolds of the configuration space
    (\cref{eq:homogeneous}).
  \item Design corresponding Finsler energies defining the importance metric in this manifold
    (\cref{sub:conservative_fabrics_and_energization}).
  \item Energize all geometries with the associated Finsler energies 
    (\cref{sub:conservative_fabrics_and_energization}).
  \item Pull back the energized systems into the configuration space and sum them 
    (\cref{sub:operations_on_specs}).
  \item Force the combined system with a differentiable potential. As a composition of optimization fabrics, 
    the resulting trajectory converges towards the potential's minimum
    (\cref{sub:optimization_fabrics}).
\end{enumerate}

This procedure has been initial presented in
\cite{Ratliff2020} and is used throughout this thesis to
generate trajectory with \ac{fabrics}. In \cref{cha:tro_23},
we adopt the procedure to our generalization to dynamic
settings.

\subsection{Experimental results fabrics}%
\label{sub:experimental_results_fabrics}

The theory explained above was tested on several simple kinematic chains
in~\cite{Ratliff2020,Ratliff2021}. As fabrics design motion as a summation of several
differential equations, each representing a specific constraint to the motion, it is
possible to sequentially design motion~\cite{Ratliff2020}. This procedure allows to
carefully tune individual components without harming the others. The application to a
planar arm in a goal-reaching setup was successfully tested in~\cite{Ratliff2020}. Here,
the authors illustrated how the resulting motion can be modified arbitrarily by the user
by adding additional constraints or preferences.

Although important concepts and findings on \ac{fabrics} were summarized in this
section, we refer to~\cite{Ratliff2020} for a more in-depth presentation of optimization
fabrics. In the following, we generalize the framework of optimization fabrics to dynamic settings.
